\documentclass[conference,9pt]{IEEEtran}

\usepackage{bm} 
\usepackage{siunitx}
\robustify\bfseries
\usepackage{booktabs}
\usepackage[capitalize,nameinlink]{cleveref}
\crefname{equation}{}{}
\Crefname{equation}{Equation}{Equations}

\title{ECSE541 Initial Report}

\author{\IEEEauthorblockN{Ryan Edrissi, Yuchen Niu, Peng Peng, Trevor Piltch, Marcus Porfilio}
\IEEEauthorblockA{Department of Electrical and Computer Engineering, McGill University\\
845 Sherbrooke Street West, Montréal H3A 0G4}}


\begin{document}

\maketitle

\section{Introduction}
\label{sec:intro}
This is the intro.


\section{Related Work}

Hill and Marty \cite{4563876} extended Amdahl's law to multicore chips with a fixed area budget. Their analysis assumes that a core's sequential performance grows roughly with the square root of the resources devoted to it. Under this assumption, spending resources on a few larger cores speeds up the sequential part of a program, while spending them on many small cores speeds up the parallel part. However, their analysis treats cores as abstract units of area and performance, and does not consider power. Kumar et al. \cite{1253185} approached the same trade-off problem by combining cores of different sizes that share one instruction set, allowing each workload to run on the core type that fits it best and reducing power.

Measuring this trade-off for real core designs requires detailed simulations. gem5 \cite{10.1145/2024716.2024718} is a cycle-level simulator that provides both in-order and out-of-order CPU models. Endo et al. \cite{6893220} used gem5 to model an ARM Cortex-A8 (in-order) and Cortex-A9 (out-of-order) and validated them by comparing simulated results and measurements from hardware. Their work focused on the accuracy of the simulations, and not on comparing the two cores. To estimate the power and area of these processors, McPAT \cite{5375438} is used to model the processor's microarchitecture. Given activity statistics from the simulator, McPAT reports chip area, timing, and dynamic and leakage power. This work rebuilds the A8 and A9 models from Endo et al., validates them against the published results, and uses McPAT to compare configurations of comparable area so the trade-off can be measured on standard benchmarks.

\section{Methodology}

We use gem5 and McPAT to investigate how core type and core count affect performance and power across workloads at comparable chip areas. We first simulate multicore configurations with varying numbers of Cortex-A8 or Cortex-A9 cores in gem5, recording the simulated execution time and runtime activity statistics for each benchmark. We then pass the processor configuration and activity statistics to McPAT to estimate chip area and runtime power. Finally, we pair Cortex-A8 and Cortex-A9 configurations with similar estimated areas and compare their execution time and power for each workload.

\subsection{Performance Simulation with gem5}

We configure an in-order core model based on the Cortex-A8 and an out-of-order core model based on the Cortex-A9 in gem5. To validate these models, we reproduce the benchmark runs described in \textcolor{red}{[TODO]}, compare our simulation results with the corresponding published results, and report the relative error.

For the multicore experiments, we construct candidate processor configurations by varying the core type and core count. For each benchmark, we use the same input and amount of work across configurations. We use simulated execution time as the performance metric and collect runtime activity statistics, such as instruction counts and cache accesses, for estimating runtime dynamic power in McPAT.

\subsection{Area and Power Estimation with McPAT}

For each processor configuration and benchmark combination, we generate a McPAT XML input file. We map parameters from the gem5 configuration, such as core count, clock frequency, and cache capacity, to their corresponding XML fields and specify additional modeling parameters, such as the process technology node. We also populate the XML file with runtime activity statistics from the corresponding simulation. McPAT uses the hardware configuration to estimate chip area and combines it with the activity statistics to estimate runtime power for the workload.

We use McPAT's area estimates to pair Cortex-A8 and Cortex-A9 configurations with similar chip areas and report the estimated area difference for each pair. For each workload, we compare the simulated execution time and estimated power of the paired configurations. This comparison shows which workloads run faster on configurations with more Cortex-A8 cores, which run faster on configurations with fewer Cortex-A9 cores, and how their power consumption differs at comparable chip areas.

\section{Experimental Setup}

\section{Conclusion}

\bibliographystyle{IEEEtran}
\bibliography{refs}

\end{document}

